%% ============================================================================
%%  lafrox-portada.tex -- Portada corporativa LafroX, variante "banda diagonal"
%%
%%  Se carga automaticamente desde lafrox.cls. Define:
%%      \portadalafrox      portada completa a sangre (folio 1)
%%      \contraportadalafrox  cierre opcional con el equipo y la declaracion
%%
%%  Toda la geometria esta parametrizada al final del archivo: para mover la
%%  diagonal basta cambiar \LFXbandaAlto, \LFXbandaCimaX y \LFXbandaBaseX.
%% ============================================================================

% ---- Parametros de la banda (fracciones del papel) -------------------------
\def\LFXbandaAlto{0.575}     % altura de la banda
\def\LFXbandaCimaX{0.34}     % ancho de la banda en el borde superior
\def\LFXbandaBaseX{0.82}     % ancho de la banda en su borde inferior

% ---- Utilidades de posicionamiento ----------------------------------------
%  \LFXen{fraccion x}{fraccion y}  -> coordenada respecto del vertice sup. izq.
%  Las llaves y el * son indispensables: sin ellas, una expresion como
%  0.34+0.095 se lee como "0.34 pt mas 0.095 anchos de papel".
\newcommand{\LFXen}[2]{%
  ([xshift={(#1)*\paperwidth},yshift={-(#2)*\paperheight}]current page.north west)}

% ---- Reticula horizontal de la portada ------------------------------------
%  \LFXmi y \LFXmd son las fracciones de papel que corresponden a los
%  margenes REALES del bloque de texto. Se calculan, no se escriben a mano:
%  asi el titulo, la reticula de datos y el pie academico de la portada caen
%  exactamente sobre los mismos margenes que cualquier pagina de contenido, y
%  siguen cayendo ahi si manana se cambia \geometry.
%
%  La banda diagonal NO usa estas fracciones: es arte a sangre y va de borde
%  a borde a proposito.
%  \LFXmc es el inicio de la segunda columna de la reticula de datos, y
%  \LFXcolportada su ancho: dos columnas de 0,46 del bloque de texto con un
%  canal de 0,08 entre ellas. Antes eran 0,375 del ANCHO DE PAPEL cada una,
%  que sumadas se salian del margen derecho.
\newcommand{\LFXcolportada}{0.46\textwidth}
\newcommand{\LFXcalcularmargenes}{%
  \pgfmathsetmacro{\LFXmi}{(1in+\oddsidemargin+\hoffset)/\paperwidth}%
  \pgfmathsetmacro{\LFXmd}{(1in+\oddsidemargin+\hoffset+\textwidth)/\paperwidth}%
  \pgfmathsetmacro{\LFXmc}{(1in+\oddsidemargin+\hoffset+0.54*\textwidth)/\paperwidth}%
}

% ---- Marca: archivo de logo si existe, monograma tipografico si no --------
\newcommand{\LFXmarcaportada}{%
  \ifx\LFXlogo\@empty
    \begin{tikzpicture}[baseline=0pt]
      \node[anchor=west,inner sep=0pt] (mono) at (0,0)
        {\lafroxmonograma[11mm]};
      \node[anchor=west,inner sep=0pt,align=left] at ([xshift=3.4mm]mono.east)
        {\sffamily\bfseries\fontsize{17pt}{17pt}\selectfont
         \color{lafroxVioleta}\LFXempresa};
    \end{tikzpicture}%
  \else
    \includegraphics[height=13mm]{\LFXlogo}%
  \fi}

% ============================================================================
%  \portadalafrox
% ============================================================================
\newcommand{\portadalafrox}{%
  \begingroup
  \setcounter{page}{1}%
  \thispagestyle{lfxportada}%
  \LFXcalcularmargenes
  \null
  \begin{tikzpicture}[remember picture,overlay]

    % --- Fondo -------------------------------------------------------------
    \fill[lafroxBg]
      (current page.south west) rectangle (current page.north east);

    % --- Diagonales de profundidad (detras de la banda principal) ----------
    \fill[lafroxAcento]
      \LFXen{\LFXbandaCimaX+0.105}{0}      --
      \LFXen{\LFXbandaCimaX+0.135}{0}      --
      \LFXen{\LFXbandaBaseX+0.135}{\LFXbandaAlto} --
      \LFXen{\LFXbandaBaseX+0.105}{\LFXbandaAlto} -- cycle;
    \fill[lafroxAzul]
      \LFXen{\LFXbandaCimaX}{0}            --
      \LFXen{\LFXbandaCimaX+0.095}{0}      --
      \LFXen{\LFXbandaBaseX+0.095}{\LFXbandaAlto} --
      \LFXen{\LFXbandaBaseX}{\LFXbandaAlto} -- cycle;

    % --- Banda principal ---------------------------------------------------
    \fill[lafroxVioleta]
      \LFXen{0}{0}                          --
      \LFXen{\LFXbandaCimaX}{0}             --
      \LFXen{\LFXbandaBaseX}{\LFXbandaAlto} --
      \LFXen{0}{\LFXbandaAlto}              -- cycle;

    % --- Marca, arriba a la derecha sobre el fondo claro -------------------
    \node[anchor=north east,inner sep=0pt] at \LFXen{\LFXmd}{0.072}
      {\LFXmarcaportada};
    \node[anchor=north east,inner sep=0pt] at \LFXen{\LFXmd}{0.140}
      {\sffamily\fontsize{9pt}{11pt}\selectfont\color{lafroxGris}\LFXcaso};

    % --- Titulo en negativo, dentro de la banda ----------------------------
    %  Dos nodos y no uno con \\: dentro de un nodo TikZ el salto de linea
    %  reinicia el tamano de fuente y la segunda linea sale minuscula.
    \node[anchor=north west,inner sep=0pt] at \LFXen{\LFXmi}{0.250}
      {\sffamily\bfseries\fontsize{9.5pt}{12pt}\selectfont
       \color{lafroxAcento}\LFXlicitacion};

    \node[anchor=north west,inner sep=0pt] at \LFXen{\LFXmi}{0.295}
      {\sffamily\bfseries\fontsize{38pt}{38pt}\selectfont\color{white}
       PROPUESTA};
    \node[anchor=north west,inner sep=0pt] at \LFXen{\LFXmi}{0.360}
      {\sffamily\bfseries\fontsize{38pt}{38pt}\selectfont\color{white}
       TÉCNICA};

    \draw[lafroxAcento,line width=3pt]
      \LFXen{\LFXmi}{0.455} -- \LFXen{\LFXmi+0.13}{0.455};

    \node[anchor=north west,inner sep=0pt] at \LFXen{\LFXmi}{0.478}
      {\sffamily\fontsize{12.5pt}{16pt}\selectfont\color{white}
       Presentada a \textbf{\LFXcliente}};

    % --- Identificacion del documento, bajo la banda -----------------------
    \node[anchor=north west,inner sep=0pt,align=left] at \LFXen{\LFXmi}{0.625}
      {\sffamily\bfseries\fontsize{23pt}{25pt}\selectfont
       \color{lafroxVioleta}\LFXsubtitulo};
    \node[anchor=north west,inner sep=0pt,align=left] at \LFXen{\LFXmi}{0.685}
      {\sffamily\fontsize{9.5pt}{13pt}\selectfont\color{lafroxGris}
       \LFXalcance\ \textperiodcentered\ \LFXformulario};

    \draw[lafroxLinea,line width=0.7pt]
      \LFXen{\LFXmi}{0.735} -- \LFXen{\LFXmd}{0.735};

    % --- Retícula de datos de la oferta ------------------------------------
    \node[anchor=north west,inner sep=0pt] at \LFXen{\LFXmi}{0.760}
      {\begin{minipage}{\LFXcolportada}
         \sffamily\fontsize{8.5pt}{13pt}\selectfont
         \color{lafroxGris}\raggedright
         \textcolor{lafroxAzul}{\bfseries Proponente}\\
         \LFXempresa\ \textperiodcentered\ RUT \LFXrut\\[5pt]
         \textcolor{lafroxAzul}{\bfseries Representante}\\
         \LFXrepresentante\ \textperiodcentered\ \LFXcargo
       \end{minipage}};

    \node[anchor=north west,inner sep=0pt] at \LFXen{\LFXmc}{0.760}
      {\begin{minipage}{\LFXcolportada}
         \sffamily\fontsize{8.5pt}{13pt}\selectfont
         \color{lafroxGris}\raggedright
         \textcolor{lafroxAzul}{\bfseries Fecha de presentación}\\
         \LFXfecha\ \textperiodcentered\ versión \LFXversion\\[5pt]
         \textcolor{lafroxAzul}{\bfseries Firma del representante}\\
         \textcolor{lafroxLinea}{\rule{4.2cm}{0.5pt}}
       \end{minipage}};

    % --- Pie academico -----------------------------------------------------
    \draw[lafroxVioleta,line width=1.1pt]
      \LFXen{\LFXmi}{0.900} -- \LFXen{\LFXmd}{0.900};
    \node[anchor=north west,inner sep=0pt,align=left] at \LFXen{\LFXmi}{0.912}
      {\sffamily\fontsize{7.5pt}{10.5pt}\selectfont\color{lafroxGris}
       \LFXasignatura\ \textperiodcentered\ \LFXunidad\\
       Profesor: \LFXprofesor};

  \end{tikzpicture}
  \clearpage
  \endgroup}

% ============================================================================
%  \contraportadalafrox -- cierre institucional (opcional)
%  Uso: \contraportadalafrox{tabla de equipo o texto libre}
% ============================================================================
\newcommand{\contraportadalafrox}[1]{%
  \clearpage
  \thispagestyle{lfxportada}%
  \null
  \begin{tikzpicture}[remember picture,overlay]
    \fill[lafroxVioleta]
      (current page.north west) rectangle
      ([yshift=-0.30\paperheight]current page.north east);
    \fill[lafroxAcento]
      ([yshift=-0.30\paperheight]current page.north west) rectangle
      ([yshift=-0.305\paperheight]current page.north east);
    \node[anchor=north west,inner sep=0pt,align=left] at \LFXen{\LFXmi}{0.115}
      {\sffamily\bfseries\fontsize{26pt}{29pt}\selectfont\color{white}
       \LFXempresa};
    \node[anchor=north west,inner sep=0pt,align=left] at \LFXen{\LFXmi}{0.200}
      {\sffamily\fontsize{10.5pt}{14pt}\selectfont\color{lafroxClaro}
       \LFXdocumento\ \textperiodcentered\ \LFXsubtitulo\\
       \LFXlicitacion\ \textperiodcentered\ \LFXcaso};
  \end{tikzpicture}
  \vspace*{0.24\paperheight}
  #1
  \clearpage}
